\section{Limitations}
\begin{itemize}\itemsep0pt
\item Sector labels are a 2026 snapshot, not point-in-time; names delisted earlier are \texttt{Unclassified}.
\item No shares-outstanding data: size is proxied by traded value; share turnover is not tested.
\item No corporate-action export: splits and bonuses are inferred by ratio matching with F\&O lot/open-interest confirmation; dividends are missing (price returns; benchmark is the price index); rights issues are captured only when their ratio is near a common one; an F\&O demerger without an open-interest jump would be kept as a fake move. Counts and the 50 most liquid inferred actions are in the data-quality report.
\item Protocol incident (disclosed): during a plumbing rehearsal on partial development data a label shuffle silently failed, and the researcher saw the combiner comparison, one signal's IC and the turnover level before the official run. Its 33 trials are in the ledger and in $N_{\mathrm{all}}$; no pre-registered choice was changed (\texttt{DECISIONS.md} D-028).
\item True delisting returns are unknown; a held name carries flat for up to 10 sessions and is then written off at a 0\% or 30\% haircut.
\item Statutory rates were spot-checked from public sources; the 2024-10-01 futures STT date and the pre-2016 delivery history are not independently confirmed (no reported number depends on pre-2016 rates). Brokerage is assumed; futures margin is a flat 20\%; lot rounding is ignored.
\item The impact model is a square-root law with $Y=1$, not fitted to Indian execution data. P2 uses cash-market ADV and volatility for the futures leg.
\item The risk-free rate is a constant 6.5\% fallback (no RBI T-bill series supplied).
\item The trial count is a lower bound: choosing these eight signals from the literature was itself a selection the ledger cannot see. The paper's own deflated-Sharpe worked example was not reproduced (the PSR example was).
\item The researcher lived through 2024--2026 markets; the lock prevents numeric peeking, not general knowledge.
\item One market, about eight years of walk-forward plus two of holdout.
\end{itemize}
